
\justifying 
Project planning is a procedural step in project management. It is the practice of initiating,
planning, executing, controlling, and closing the work team to achieve specific goals. Project
planning and management are important because it ensures that the right people do the right
things, at the right time. It also ensures the proper project life cycle.

The organization of the Chapter is as follows . The Section 2.1 shows the Feasibility
Study of the project. Risk Analysis of the project is represented in Section 2.2. Project
Scheduling is described in Section 2.3. The Section 2.4  describes the effort allocation, Cost Estimation given in Section 2.5 respectively. Summary is mentioned in Section 2.6.


\section{Feasibility Study}

A feasibility study is an important step in software development that evaluates the practicality and success potential of a project. In this project, namely the Virtual Lab Assistant for DBMS practical examination, the feasibility study helps in identifying whether the system can be developed and implemented effectively.

The study focuses on analyzing different aspects such as technical feasibility, operational feasibility, and economic feasibility. It helps in understanding the strengths, weaknesses, opportunities, and possible risks associated with the system. This ensures that the project is realistic and achievable within the available resources and time.

The system is designed as a web-based platform using modern technologies such as Firebase, HTML, CSS, and JavaScript. It provides features like student authentication, practical exam execution, performance tracking, and secure session handling. The feasibility study evaluates whether these features can be efficiently developed and maintained.



\subsection{Technical Feasibility}
The technical feasibility of the proposed VR-Lab Assistant system is evaluated based on the system requirements and the technologies used for development. The objective is to determine whether the available technical resources are sufficient to successfully design and implement the system.

The project is developed using HTML, CSS, and JavaScript for the frontend, along with Firebase (Authentication and Firestore Database) for backend services. Firebase provides real-time database support, secure authentication, and cloud-based storage, eliminating the need for complex server-side setup


\subsection{Operational Feasibility}

Operational feasibility evaluates how effectively the VR-Lab Assistant system will function within the existing educational environment.

The system is designed to assist students in performing virtual practicals and conducting online examinations in a structured and interactive way. It provides features such as:
\begin{enumerate}
\item Student login and authentication
\item Admin control panel
\item Practical execution with real-time query evaluation
\item Performance tracking and progress visualization
\end{enumerate}
\subsection{Economic Feasibility}
Economic feasibility determines whether the project is cost-effective and beneficial compared to traditional methods.

The VR-Lab Assistant system is highly economical as it uses free or low-cost technologies such as Firebase and web-based tools. There is no need for expensive hardware, licensed software, or dedicated servers.



\section{Risk Analysis}
• The Virtual Lab Assistant project faces potential risks in technical, security, and
operational areas that could affect its smooth execution and reliability. Key concerns
include dependency on APIs, database performance, system security, and user
adoption. However, with planned mitigation strategies like fallback systems, testing,
monitoring, and strict schedule adherence, these risks can be effectively managed to
ensure stable and secure project delivery.

\item  Technical Risks


• Possible query evaluation inaccuracies affecting grading credibility.

• Firebase scaling or outage risks causing downtime or performance issues.

• Security and Integrity Risks

• Attempts to bypass exam security features like tab restriction or copy-paste blocking.

• Potential data breaches exposing student or exam information.

• Operational and Usability Risks

• Resistance to platform adoption by faculty or students.


• Delays in project completion due to the tight four-week schedule. 

\item 
\item 
\item 

\section{Project Scheduling}
Project scheduling is an essential part of project planning that defines the timeline and sequence of activities required for successful completion of the project. It helps in organizing tasks, allocating time efficiently, and ensuring that the project progresses in a structured manner. A well-defined schedule allows developers to track progress, identify delays, and manage resources effectively.

 
\raggedright 
\subsection{Week-by-Week Breakdown}

\begin{itemize}

\item \textbf{Week 1:} Focuses primarily on Planning / Requirements and initiating the Analysis \& Design phase.

\item \textbf{Week 2:} The Analysis \& Design phase continues, and the bulk of the Implementation / Development work begins.

\item \textbf{Week 3:} Is the most intensive week, featuring major activities in Implementation / Development, the commencement of Testing, and the start of Review / Evaluation.

\item \textbf{Week 4:} Concludes the project with the final activities for Implementation / Development and Testing. It also includes the continuation of Review / Evaluation and the final stage of Deployment / Integration.

\end{itemize}
\end{figure}
\newpage

\justifying 
\section{Effort Allocation}
Effort allocation refers to the amount of time a task should take from an assignee’s day. A project involves teamwork, where development is the result of combined efforts from the team. The project is divided into modules, with a number of modules allocated to team members. After completion of each module, it will be linked to other modules to form a complete project. Effort allocation should serve as a guideline only. The characteristics of each project will determine the distribution of effort. The effort allocation table for the proposed system is shown in the following table. Table 2.1 shows Effort Allocation for Project.
%\centering

%\text{Table 2.2 : Effort Allocation }\vspace{0.5cm}
%\index{ Effort Allocation for Project}
\begin{table}[H]
  %\centering
  % Your table content here
  \caption{Effort Allocation for Project}\index{ Effort Allocation for Project}
\end{table}


   \begin{tabular}{|p{0.6cm}|p{6.7cm}|p{1.8 cm}|p{1.9 cm}|p{1.8cm}|p{1.8cm}|}
       
        \hline
        Sr. no & Module & Pradnya & Chhayshree & Gauri & Saniya \\ \hline
        1 & Gathering Of Information & \checkmark& & \checkmark & \\ \hline
        2 & Planing/Requirement Analysis & \checkmark  & \checkmark &  & \checkmark \\ \hline
        3 & Selection of life cycle &  & \checkmark & \checkmark & \checkmark \\ \hline
        4 & Planning and management & \checkmark & \checkmark &  & \\ \hline
        5 & Analysis \& Design UML  & \checkmark & \checkmark & \checkmark & \checkmark \\ \hline
        
   \end{tabular}]

%\begin{table}[H]
%  \centering
  % Your table content here
 % \caption{Effort Allocation for Project}
%\end{table}

\newline
\section{Cost Estimation}


The Virtual Lab Assistant project is highly cost-effective, leveraging free and freemium tools to minimize expenses. With HTML, CSS, and JavaScript for the frontend and Firebase (using Google’s free Spark Plan) for the backend, initial development costs remain minimal. The only significant variable expense arises from the OpenAI API, which operates on a pay-as-you-go model based on hint usage. Overall, the project maintains a low initial cost structure, with expenses scaling only as user activity increases.

The cost estimation model used in this project follows a \textbf{Zero-Based Budgeting approach} combined with a \textbf{Lean/Startup Cost Structure}. In this approach, development begins with an estimated cost of ₹0 for most components such as frontend development, backend development, database, testing, and maintenance, and expenses are added only where absolutely necessary (such as domain name, contingency, and API usage).

\subsubsection{Key Cost Factors}

\begin{itemize}

\item \textbf{Frontend:} Built using free technologies such as HTML, CSS, and JavaScript, resulting in no licensing cost.

\item \textbf{Backend \& Database:} Hosted on Firebase Freemium (Spark Plan), where costs scale based on usage.


\item \textbf{Operational Costs:} Includes Firebase reads/writes, hosting, and storage, all of which are usage-dependent.

\item \textbf{Cost Optimization:} A preloaded hints database is used as an alternative to reduce dependency on the OpenAI API and minimize costs.

\end{itemize}

The project provides a low-cost and scalable solution with minimal upfront investment and manageable variable expenses, making it suitable for educational environments.

\newline


\section{Summary}
Project Scheduling and Effort Allocation
The project follows an intensive four-week schedule to manage the workload from conception
to deployment. Effort allocation is distributed across six phases: 



